%=========== ΛΥΣΗ - 88ο ΘΕΜΑ Γ ===========

\begin{thema}{Γ}

\begin{center}
\begin{tikzpicture}[scale=0.65]
  \draw[help lines, gray!30] (-1,-4) grid (7,3);
  \draw[-latex] (-1.3,0) -- (7.3,0) node[right] {$x$};
  \draw[-latex] (0,-4.3) -- (0,3.3) node[above] {$y$};
  \foreach \x in {1,2,3,4,5,6}
    \draw (\x,0.1) -- (\x,-0.1) node[below,font=\footnotesize] {$\x$};
  \foreach \y in {-3,-2,-1,1,2}
    \draw (0.1,\y) -- (-0.1,\y) node[left,font=\footnotesize] {$\y$};

  \draw[very thick, secondcolor, smooth, tension=0.5]
    plot coordinates {(0,2) (1,1) (2,0) (3,-2) (4,-3) (5,-2) (6,-1)};

  \filldraw[secondcolor] (0,2) circle (2pt) node[above right,font=\footnotesize] {$(0,2)$};
  \filldraw[secondcolor] (2,0) circle (2pt) node[above right,font=\footnotesize] {$(2,0)$};
  \filldraw[secondcolor] (4,-3) circle (2pt) node[below right,font=\footnotesize] {$(4,-3)$};
  \node[secondcolor] at (6,1.5) {$C_{f''}$};
\end{tikzpicture}
\end{center}

\begin{erwthma}
\item Από το διάγραμμα της $f''$ παρατηρούμε ότι:
\[
f''(x)>0,\quad x\in[0,2)
\]
και
\[
f''(x)<0,\quad x\in(2,6].
\]

Επομένως, η $f$ είναι κυρτή στο $[0,2]$ και κοίλη στο $[2,6]$.

Στο $x=2$ η $f''$ μηδενίζεται και αλλάζει πρόσημο, από θετική σε αρνητική.
Άρα η $C_f$ παρουσιάζει σημείο καμπής στο σημείο
\[
\bigl(2,f(2)\bigr).
\]

\item Επειδή
\[
f''(x)>0,\quad x\in[0,2),
\]
η $f'$ είναι γνησίως αύξουσα στο $[0,2]$.

Επίσης, επειδή
\[
f''(x)<0,\quad x\in(2,6],
\]
η $f'$ είναι γνησίως φθίνουσα στο $[2,6]$.

Ακόμη, δίνεται ότι $f'(0)=0$. Επειδή η $f'$ είναι γνησίως αύξουσα στο $[0,2]$, για κάθε $x\in(0,2]$ ισχύει
\[
f'(x)>f'(0)=0.
\]
Άρα η $f$ είναι γνησίως αύξουσα στο $[0,2]$.

Στο διάστημα $(2,6)$ η $f'$ είναι γνησίως φθίνουσα. Επομένως, η μονοτονία της $f$ σε αυτό το διάστημα θα εξαρτηθεί από το αν η $f'$ μηδενίζεται μετά το $2$. Αυτό θα εξεταστεί στο τελευταίο ερώτημα.

\item Από το Θεμελιώδες Θεώρημα του Ολοκληρωτικού Λογισμού έχουμε
\[
f'(2)-f'(0)=\int_0^2 f''(x)\,\d x.
\]
Επειδή $f'(0)=0$, προκύπτει
\[
f'(2)=\int_0^2 f''(x)\,\d x.
\]

Στο διάστημα $[0,2]$ η γραφική παράσταση της $f''$ βρίσκεται πάνω από τον άξονα $x'x$ και δεν ταυτίζεται με αυτόν. Άρα
\[
\int_0^2 f''(x)\,\d x>0.
\]
Επομένως,
\[
f'(2)>0.
\]

\item Για να εξετάσουμε αν η $f$ παρουσιάζει τοπικό μέγιστο για κάποιο $\xi>2$, μελετούμε το πρόσημο της $f'$ στο $(2,6)$.

Από το προηγούμενο ερώτημα έχουμε
\[
f'(2)>0.
\]
Επιπλέον,
\[
f'(6)-f'(2)=\int_2^6 f''(x)\,\d x,
\]
οπότε
\[
f'(6)=f'(2)+\int_2^6 f''(x)\,\d x.
\]

Από το διάγραμμα, το θετικό εμβαδό ανάμεσα στην $C_{f''}$ και τον άξονα $x'x$ στο διάστημα $[0,2]$ είναι μικρότερο από $4$, αφού η $f''$ βρίσκεται κάτω από την ευθεία $y=2$:
\[
0<f'(2)=\int_0^2 f''(x)\,\d x<4.
\]

Αντίστοιχα, στο διάστημα $[2,6]$ η $C_{f''}$ βρίσκεται κάτω από τον άξονα $x'x$. Μάλιστα, από το σχήμα, στο διάστημα $[3,5]$ η καμπύλη βρίσκεται κάτω από την ευθεία $y=-2$ και δεν ταυτίζεται με αυτήν. Άρα
\[
\int_3^5 f''(x)\,\d x<-4.
\]
Επειδή και στα διαστήματα $[2,3]$, $[5,6]$ ισχύει $f''(x)<0$, έχουμε τελικά
\[
\int_2^6 f''(x)\,\d x<-4.
\]

Άρα
\[
f'(6)=f'(2)+\int_2^6 f''(x)\,\d x<4-4=0.
\]
Επομένως,
\[
f'(2)>0
\quad\text{και}\quad
f'(6)<0.
\]

Η $f'$ είναι συνεχής στο $[2,6]$, αφού η $f$ είναι δύο φορές παραγωγίσιμη. Άρα, από το Θεώρημα \eng{Bolzano}, υπάρχει $\xi\in(2,6)$ τέτοιο ώστε
\[
f'(\xi)=0.
\]

Επειδή η $f'$ είναι γνησίως φθίνουσα στο $[2,6]$, το $\xi$ είναι μοναδικό. Επιπλέον,
\[
f'(x)>0,\quad x\in(2,\xi)
\]
και
\[
f'(x)<0,\quad x\in(\xi,6).
\]

Άρα η $f$ είναι γνησίως αύξουσα στο $(2,\xi)$ και γνησίως φθίνουσα στο $(\xi,6)$. Συνεπώς, η $f$ παρουσιάζει τοπικό μέγιστο στο $x=\xi$, με $\xi>2$.
\end{erwthma}
\end{thema}

%=========== ΤΕΛΟΣ ΛΥΣΗΣ - 88ο ΘΕΜΑ Γ ===========
